\clearpage\thispagestyle{empty}
\pdfbookmark[0]{Quellen zum Brief vom 3. Januar 1807}{zach18070103-sources}
\begin{center}\textbf{Quellen zum Brief vom 3. Januar 1807}\end{center}
\noindent Diese KI-Übertragung in heutiges Deutsch folgt dem Brief Carl Friedrich Gauß’ vom 3. Januar 1807 an Franz Xaver von Zach in \href{https://gdz.sub.uni-goettingen.de/id/PPN236007467}{Gauß’ \textit{Werken}, Band~VI} (Göttingen, 1874), S.~280--281. Er erschien als Beitrag VIII in der \textit{Monatlichen Correspondenz}, Band~XV, Februar 1807, \href{https://archive.org/details/bub_gb_sKAAAAAAMAAJ/page/n157/mode/1up}{S.~152}--\href{https://archive.org/details/bub_gb_sKAAAAAAMAAJ/page/n162/mode/1up}{157}. Von Zachs Anmerkung zur verspäteten Post steht auf \href{https://archive.org/details/bub_gb_sKAAAAAAMAAJ/page/n157/mode/1up}{S.~152}. Die Ephemeride beginnt auf \href{https://archive.org/details/bub_gb_sKAAAAAAMAAJ/page/n159/mode/1up}{S.~154}. Dort berichtet von Zach in einer eigenen Anmerkung, dass Hardings Karten für Ceres, Pallas und Juno bereits beim Kupferstecher waren.
\par\medskip
Die beiden Verweise auf Pasquichs Beobachtungen führen zu \href{https://www.digitale-sammlungen.de/en/view/bsb10538605?page=200}{Band~XIII, S.~192} und \href{https://www.digitale-sammlungen.de/en/view/bsb10538606?page=102}{Band~XIV, S.~92}. Im folgenden Beitrag, \href{https://archive.org/details/bub_gb_sKAAAAAAMAAJ/page/n163/mode/1up}{Band~XV, S.~158}--\href{https://archive.org/details/bub_gb_sKAAAAAAMAAJ/page/n164/mode/1up}{159}, übersendet Oriani die Mailänder Beobachtungen und nennt Carlini als Urheber der Beobachtungen und Berechnungen. In der Beschreibung der Ephemeride in eckigen Klammern steht in der Werkausgabe \emph{Delimalen}. Gemeint sind Dezimalstellen.
\par\medskip
Das Februarheft stammt aus dem von Google digitalisierten Exemplar der New York Public Library im Internet Archive. Die Hefte von 1806 stammen aus der Bayerischen Staatsbibliothek. Übertragung und Satz wurden mit KI erstellt.
\clearpage
